\documentclass[11pt]{article}
\usepackage{amsmath, amssymb}
\usepackage{geometry}
\usepackage{enumitem}
\usepackage{booktabs}
\usepackage{graphicx}
\geometry{margin=1in}

\title{STAT 17 Week 8 \\Confidence Intervals}
\author{Xiao Zhang}
\date{February 2026}

\begin{document}

\maketitle

\section*{1. Statistical Inference Overview}

Statistical inference allows us to use sample data to estimate population parameters.

\textbf{Target Parameter}
\begin{itemize}
    \item Quantitative data $\rightarrow$ population mean $\mu$
    \item Categorical data $\rightarrow$ population proportion $p$
\end{itemize}

\textbf{Point Estimate}
\begin{itemize}
    \item $\bar{x}$ estimates $\mu$
    \item $\hat{p}$ estimates $p$
\end{itemize}

\textbf{Confidence Interval (CI)}
\[
\text{Point Estimate} \pm (\text{Critical Value})(\text{Standard Error})
\]

\section*{2. Confidence Level Interpretation}

\textbf{Correct Interpretation:}

In the long run, about $(1-\alpha)\%$ of all possible samples will produce confidence intervals that contain the true parameter.

\textbf{Important:} Confidence is NOT probability.

We do NOT say:
\begin{itemize}
    \item ``There is a 95\% probability the parameter is inside.''
    \item ``95\% of data values fall inside.''
\end{itemize}

\section*{3. Confidence Interval for a Population Proportion}

\subsection*{When to Use}
\begin{itemize}
    \item Data are Yes/No (success/failure)
    \item Interested in population proportion $p$
\end{itemize}

\subsection*{Conditions}
\begin{enumerate}
    \item Random sample
    \item $n \le 0.05N$
    \item $n\hat{p}(1-\hat{p}) \ge 10$
\end{enumerate}

\subsection*{Formula}
\[
\hat{p} \pm z_{\alpha/2} \sqrt{\frac{\hat{p}(1-\hat{p})}{n}}
\]

\textbf{Critical values (Z-table):}

\begin{center}
\begin{tabular}{cc}
\toprule
Confidence Level & $z_{\alpha/2}$ \\
\midrule
90\% & 1.645 \\
95\% & 1.96 \\
99\% & 2.575 \\
\bottomrule
\end{tabular}
\end{center}

\section*{4. Confidence Interval for a Population Mean}

\subsection*{Case 1: $\sigma$ Known (or large $n$)}

\subsubsection*{When to Use}
\begin{itemize}
    \item Interested in $\mu$
    \item Population standard deviation $\sigma$ known
    \item OR large sample size ($n \ge 30$)
\end{itemize}

\subsubsection*{Formula}
\[
\bar{x} \pm z_{\alpha/2} \frac{\sigma}{\sqrt{n}}
\]

\subsection*{Case 2: $\sigma$ Unknown (Use $s$)}

\subsubsection*{When to Use}
\begin{itemize}
    \item Interested in $\mu$
    \item $\sigma$ unknown (almost always)
    \item Small sample ($n < 30$)
\end{itemize}

\subsubsection*{Conditions}
\begin{enumerate}
    \item Random sample
    \item $n \le 0.05N$
    \item Population approximately normal OR $n \ge 30$
\end{enumerate}

\subsubsection*{Formula}
\[
\bar{x} \pm t_{\alpha/2, df} \frac{s}{\sqrt{n}}
\]
\[
df = n - 1
\]

\section*{5. Z vs T Decision Rule}

\begin{center}
\begin{tabular}{cccc}
\toprule
Parameter & $\sigma$ Known? & Sample Size & Use \\
\midrule
Proportion & N/A & Any & Z \\
Mean & Yes & Any & Z \\
Mean & No & $n < 30$ & T \\
Mean & No & $n \ge 30$ & Usually T \\
\bottomrule
\end{tabular}
\end{center}

\textbf{Key Rule:}
\begin{itemize}
    \item Proportion $\rightarrow$ Always Z
    \item Mean $\rightarrow$ Check if $\sigma$ known
\end{itemize}

\section*{6. Margin of Error}

\[
\text{ME} = (\text{Critical Value})(\text{Standard Error})
\]

\textbf{Effects on Width:}
\begin{itemize}
    \item Increase confidence level $\rightarrow$ wider interval
    \item Increase sample size $\rightarrow$ narrower interval
    \item If $n$ quadruples $\rightarrow$ ME is cut in half
\end{itemize}

\section*{7. Sample Size Formulas}

\subsection*{For Proportion}
\[
n = \hat{p}(1-\hat{p})\left(\frac{z_{\alpha/2}}{E}\right)^2
\]

If no estimate available, use $\hat{p} = 0.5$.

\subsection*{For Mean}
\[
n = \left(\frac{z_{\alpha/2}s}{E}\right)^2
\]

Always round up.

\section*{8. Interpretation Templates}

\textbf{Confidence Interval:}

We are \_\_\_\% confident that the true (parameter in context) lies between (lower bound) and (upper bound).

\textbf{Confidence Level:}

In the long run, about \_\_\_\% of all possible samples will produce intervals that contain the true parameter.

\section*{9. Common Exam Traps}

\begin{itemize}
    \item Mixing up parameter vs statistic
    \item Using Z instead of T for small samples
    \item Incorrect interpretation of confidence level
    \item Forgetting to check conditions
    \item Forgetting degrees of freedom ($df=n-1$)
\end{itemize}

\end{document}